\chapter{Семантики языков программирования}

\section{Язык выражений}

$$ \llb \cdot \rrb : \mathcal E \to (\Sigma \to \Z) $$
$$ \llb \cdot \rrb_s : \mathcal E \to \underbrace{\Set{\text{Int}, \text{Bool}}_\bot}_T $$

$ \Gamma \vdash e : T $ "--- judgement (суждение).

$$ \Gamma \vdash e : \text{Bool} $$
$$ \Gamma \vdash e : \text{Int} $$
$$ \Gamma \vdash e : \bot $$

$$ \llb \text{Int} \rrb = \Z, \quad \llb \text{Bool} \rrb = \Set{0, 1}, \quad \llb \bot \rrb = \bot $$
$$ \llb \Gamma = \Set{\sigma | \forall x \quad \sigma x \in \llb \Gamma x \rrb} $$

\begin{undefthm}{Наша семантика.}
	$$ \llb n \rrb = \sigma \mapsto w $$
	$$ \llb x \rrb = \sigma \mapsto \sigma x $$
	$$ \llb l \otimes r \rrb = \sigma \mapsto \llb l \rrb \oplus \llb r \rrb $$
\end{undefthm}

Такой стиль описания семантик называется \emph{денотационным}.

\begin{undefthm}{Вторая наша семантика.}
	$$ \Gamma \vdash x : \Gamma x, \quad \Gamma \vdash n : \text{Int} $$
	$$ \frac{\Gamma \vdash l : \text{Int} \quad r : \text{Int}}{\Gamma \vdash l
		\begin{subarray}{c}
			+ \\
			- \\
			*
		\end{subarray} r : \text{Int}}, \quad
		\frac{\Gamma \vdash l : \text{Int} \quad \Gamma \vdash r : \text{Int}}{\Gamma \vdash l
		\begin{subarray}{c}
			/ \\
			\%
		\end{subarray} r : \bot} $$
	$$ \frac{\Gamma \vdash l : \text{Int} \quad \Gamma \vdash r : \text{Int}}{\Gamma \vdash l
		\begin{subarray}{c}
			< \\
			\le \\
			> \\
			\ge \\
			= \\
			\ne
		\end{subarray} r : \text{Bool}}, \quad
	\frac{\Gamma \vdash l : \text{Bool} \quad \Gamma \vdash r : \text{Bool}}{\Gamma \vdash l
		\begin{subarray}{c}
			\vee \\
			\amp
		\end{subarray} r : \text{Bool}} $$
	И ещё одно правило, когда один из аргументов не определён (тогда и результат не определён).
\end{undefthm}

\begin{statement}
	$$ \forall \Gamma, e, T \quad \Gamma \vdash e : T \stackrel?\iff
	\forall \sigma \in \llb \Gamma \rrb \quad \llb e \rrb_\sigma \in \llb T \rrb $$
\end{statement}

$ \implies $ называется ``\emph{безопасность}'' (\emph{type safety}).
$ \impliedby $ называется ``\emph{полнота}'' (\emph{type completeness}).

\begin{proof}[безопасность]
	Структурная индукция.
\end{proof}

\begin{proof}[неполнота]
	Нужно доказать, что
	$$ \forall \sigma ~ \forall e \quad \llb e \rrb_\sigma = z \stackrel?\implies
	\sigma \in \llb \Gamma \rrb \quad \Gamma \vdash e : T \quad z \in \llb T \rrb $$

	Это верно не всегда:
	$$ \sigma = [x = 1, ~ y = 1] \quad \llb x \vee y \rrb_\sigma = 1 $$
	$ \sigma $ лежит в такой интерпретации:
	$$ \sigma \in \llb x : \text{Int}, ~ y : \text{Int} \rrb \vdash x \vee y : \bot $$
\end{proof}

Можно считать, что существует только один тип (и неопределённость тоже лежит в этом типе), то будет
и полнота и безопасность, но это бесполезно.
Почти для любой нетривиальной безопасной системы типов, полноты не будет (по теореме Гильберта).

Рассмотрим выражение $ e_1 / e_2 $.
Сейчас мы ему приписываем тип $ \bot $, поскольку $ e_2 $ может быть равно нулю.
Если бы мы могли доказать, что $ e_2 $ не обращается в ноль (в некоторых случаях это можно
доказать), то можно было бы ввести правило
$$ \frac{\Gamma \vdash e_1 : \text{Int} \quad \Gamma \vdash e_2 : \text{Int} \quad \forall \sigma ~
\llb e_2 \rrb_\sigma \ne 0}{\Gamma \vdash e_1/e_2 : \text{Int}} $$
Но $ \forall \sigma ~ \llb e_2 \rrb_\sigma \ne 0 $ "--- это одна неразрешимая проблема.

\section{Ещё два способа писания семантик}

\subsection[Операционная семантика большого шага]
{Операционная семантика большого шага (\textenglish{Big-Step Operational Semantics})}

Для начала надо описать \emph{отношение перехода} (\emph{transition relation}):
$$ \text{``$ \implies $''} \subseteq \Sigma \times \mathcal E \times \Z $$
Обычно используется инфиксная запись:
$$ \underset{\in \Sigma}c \xRightarrow{e \in \mathcal E} \underset{\in \Z}z $$

Определим отношение перехода для языка выражений:
$$ \sigma \xRightarrow{n} n $$
$$ \sigma \xRightarrow{x} \sigma x $$
$$ \frac{\sigma \xRightarrow{e} x \quad \sigma \xRightarrow{r} y}
{\sigma \xRightarrow{e \otimes r} x \oplus y} $$

Теперь семантика определяется так:
$$ \frac{\sigma \xRightarrow{e} z}{\llb e \rrb_\sigma = z} $$

$$ \llb (x + y) * (z + r) \rrb $$
В каком порядке надо посчитать множители?
Эта семантика нам не даёт ответа.
$$ \frac{\dfrac{\sigma \dots}{\sigma \xRightarrow{x + y} \dots} \quad
\dfrac{\sigma \xRightarrow{z} \sigma z \quad \sigma \xRightarrow{r} \sigma r}
{\sigma \xRightarrow{z + r} \dots}}
{\sigma \xRightarrow{(x + y) * (z + r)} \dots} $$

\subsection{Операционная семантика малого шага
(\texorpdfstring{\textenglish{Small-Step Operational Semantics}}{Small-Step Operational Semantics})}

$$ \text{``$ \to $''} \subseteq \Sigma \times \mathcal E \times \Z \cup
\Sigma \times \mathcal E \times \mathcal E $$

$$ \sigma \xrightarrow{z_1 \otimes z_2} z_1 \oplus z_2, \quad \sigma \xrightarrow{x} \sigma x $$

$$ \frac{\sigma \xrightarrow{e} e', \quad e \notin \Z}{\sigma \xrightarrow{e \otimes r}{e' \otimes r}} $$
$$ \frac{\sigma \xrightarrow{e} z, \quad e \notin \Z}{\sigma \xrightarrow{e \otimes r}{z \otimes r}} $$
$$ \frac{\sigma \xrightarrow{\Gamma}{z'}}{\sigma \xrightarrow{z \otimes r} z \oplus z'} $$
$$ \frac{\sigma \xrightarrow{r} r'}{\sigma \xrightarrow{z \otimes r} z \otimes r} $$

В такой семантике мы строим последовательность вывода:
$$ \sigma \xrightarrow{e} r', \quad \sigma \xrightarrow{e'} e'', \quad \dots $$

Семантики большого и малого шага эквивалентны функционально (то есть, одинаковым программам они
приписывают одинаковые результаты).
Но, несмотря на это, у них разная доказательная сила: в семантике малого шага нельзя доказать, что
программа, которая падает, не эквивалентна программе, которая зацикливается.
